% Owner: Member C (Eesha Irfan); mapping assumed from proposal order.
\chapter{High-Level and Low-Level Design}
Provide the high- and low-level design of your system in this chapter. Select the design that is appropriate for your project.\\
\textcolor{red}{NOTE: This Chapter is compulsory for both R\&D projects and Development projects.}
\section{System Overview}
Provide a general description of the software system, including its functionality and matters related to the overall system and its design (perhaps including a discussion of the basic design approach or organization). Feel free to split this discussion up into subsections (and sub-subsections, etc. ...).
\section{Design Considerations}
This section describes many issues that need to be addressed or resolved before attempting to devise a complete design solution.
\subsection{Assumptions and Dependencies}
Describe any assumptions or dependencies regarding the software and its use. These may concern such issues as:
\begin{itemize}
\item Related software or hardware 
  \item Operating systems 
  \item End-user characteristics 
  \item Possible and/or probable changes in functionality
\end{itemize}
Do not just add random generic statements, be clear and concise about what you are writing.
\subsection{General Constraints}
Describe any global limitations or constraints that have a significant impact on the design of the system's software (and describe the associated impact). Such constraints may be imposed by any of the following (the list is not exhaustive):
\begin{itemize}
\item Hardware or software environment 
\item End-user environment 
\item Availability or volatility of resources 
\item Standards compliance 
\item Interoperability requirements 
\item Interface/protocol requirements 
\item Data repository and distribution requirements 
\item Security requirements (or other such regulations) 
\item Memory and other capacity limitations 
\item Performance requirements 
\item Network communications 
\item Verification and validation requirements (testing) 
\item Other means of addressing quality goals 
\item Other requirements described in the requirements specification 
\end{itemize}

\subsection{Goals and Guidelines}
Describe any goals, guidelines, principles, or priorities which dominate or embody the design of the system's software. Such goals might be:
\begin{itemize}
\item The KISS principle ("Keep it simple stupid!") 
\item Emphasis on speed versus memory use 
\item Working, looking, or "feeling" like an existing product
\end{itemize}
For each such goal or guideline, unless it is implicitly obvious, describe the reason for its desirability. Feel free to state and describe each goal in its own sub subsection if you wish.
\subsection{ Development Methods}
Briefly describe the method or approach used for this software design. If one or more formal/published methods were adopted or adapted, then include a reference to a more detailed description of these methods. If several methods were seriously considered, then each such method should be mentioned, along with a brief explanation of why all or part of it was used or not used.
\section{System Architecture}
This section should describe both the internal architecture of the modules, as well as the external architecture of the system with other systems, if any. Diagrammatic architecture is compulsory.\\
This section should provide a high-level overview of how the functionality and responsibilities of the system were partitioned and then assigned to subsystems or components. Do not go into too much detail about the individual components themselves (there is a subsequent section for detailed component descriptions). The main purpose here is to gain a general understanding of how and why the system was decomposed, and how the individual parts work together to provide the desired functionality.\\
At the top-most level, describe the major responsibilities that the software must undertake and the various roles that the system (or portions of the system) must play. Describe how the system was broken down into its components/subsystems (identifying each top-level component/subsystem and the roles/responsibilities assigned to it). Describe how the higher-level components collaborate with each other in order to achieve the required results. Don't forget to provide some sort of rationale for choosing this particular decomposition of the system (perhaps discussing other proposed decompositions and why they were rejected). Feel free to make use of design patterns, either in describing parts of the architecture (in pattern format), or for referring to elements of the architecture that employ them.\\
If there are any diagrams, models, flowcharts, documented scenarios or use-cases of the system behavior and/or structure, they may be included here (unless you feel they are complex enough to merit being placed in the Detailed System Design section). Diagrams that describe a particular component or sub-system should be included within the particular subsection that describes that component or subsystem.\\
\textcolor{red}{Note: This section (and its subsections) really applies only to newly developed (or yet-to-be developed) portions of the system.} If there are parts of the system that already existed before this development effort began, then you only need to describe the pre-existing parts that the new parts of the system depend upon, and only in enough detail sufficient to describe the relationships and interactions between the old parts and the new parts. Pre-existing parts that are modified or enhanced need to be described only to the extent that is necessary for the reader to gain a sufficient understanding of the nature of the changes that were made.
\subsection{Subsystem Architecture}
If a particular component is one which merits a more detailed discussion than what was presented in the System Architecture section, provide that more detailed discussion in a subsection of the System Architecture section (or it may even be more appropriate to describe the component in its own design document). If necessary, describe how the component was further divided into subcomponents, and the relationships and interactions between the subcomponents (similar to what was done for top-level components in the System Architecture section).\\
If any subcomponents are also deemed to further merit discussion, then describe them in a separate subsection of this section (and in a similar fashion). Proceed to go into as many levels/subsections of discussion as needed in order for the reader to gain a high-level understanding of the entire system or subsystem (but remember to leave the gory details for the Detailed System Design section).\\
If this component is very large and/or complex, you may want to consider documenting its design in a separate document and simply including a reference to it in this section. If this is the option you choose, the design document for this component should have an organizational format that is very similar (if not identical to) this document.
\section{Architectural Strategies}
Explain all strategies that you use for designing of the Architecture. Describe any design decisions and/or strategies that affect the overall organization of the system and its higher-level structures. These strategies should provide insight into the key abstractions and mechanisms used in the system architecture. Describe the reasoning employed for each decision and/or strategy (possibly referring to previously stated design goals and principles) and how any design goals or priorities were balanced or traded-off. Such decisions might concern (but are not limited to) things like the following:
\begin{itemize}
    \item Use of a particular type of product (programming language, database, library, etc. ...) 
  \item Reuse of existing software components to implement various parts/features of the system 
  \item Future plans for extending or enhancing the software 
  \item User interface paradigms (or system input and output models) 
  \item Hardware and/or software interface paradigms 
  \item Error detection and recovery 
  \item Memory management policies 
  \item External databases and/or data storage management and persistence 
  \item Distributed data or control over a network 
  \item Generalized approaches to control 
  \item Concurrency and synchronization 
  \item Communication mechanisms 
  \item Management of other resources
\end{itemize}
\subsection{Architectural Strategies Strategy-1 name or description} 
\subsection{Architectural Strategies Strategy-2 name or description} 
Each significant strategy employed should probably be discussed in its subsection or (if it is complex enough) in a separate design document (with an appropriate reference here, of course). Make sure that when describing a design decision you also discuss any other significant alternatives that were considered and your reasons for rejecting them (as well as your reasons for accepting the alternative you finally chose). Sometimes it will be most effective to employ the "pattern format" for describing a strategy.
\section{Domain Model/Class Diagram}
In this subsection, add the Class Diagram of your system. Class diagrams represent the structure design of your system.

\section{Policies and Tactics}
Describe any design policies and/or tactics that do not have sweeping architectural implications (meaning they would not significantly affect the overall organization of the system and its high-level structures) but which nonetheless affect the details of the interface and/or implementation of various aspects of the system. Such decisions might concern (but are not limited to) things like the following:
\begin{itemize}
    \item Choice of which specific product to use (compiler, interpreter, database, library, etc. ...) 
 \item Engineering trade-offs 
 \item	Coding guidelines and conventions 
 \item	The protocol of one or more subsystems, modules, or subroutines 
 \item The choice of a particular algorithm or programming idiom (design pattern) to implement portions of the system's functionality 
 \item	Plans for ensuring requirements traceability 
 \item	Plans for testing the software 
 \item	Plans for maintaining the software 
 \item	Interfaces for end-users, software, hardware, and communications 
 \item	Hierarchical organization of the source code into its physical components (files and directories). 
 \item	How to build and/or generate the system's deliverables (how to compile, link, load, etc. ...) 
\end{itemize}
\subsection{Policy/tactic-1 name or description}
\subsection{Policy/tactic-2 name or description}
Each particular policy or set of tactics employed should probably be discussed in its own subsection, or (if it is large or complex enough) in a separate design document (with an appropriate reference here of course). Make sure that when describing a design decision that you also discuss any other significant alternatives that were considered, and your reasons for rejecting them (as well as your reasons for accepting the alternative you finally chose). For this reason, it may frequently be convenient to use one of the more popular "pattern formats" to describe a given tactic.\\
For this particular section, it may become difficult to decide whether a particular policy or set of tactics should be discussed in this section, or in the System Architecture section, or in the Detailed System Design section for the appropriate component. You will have to use your own "best" judgment to decide this. There will usually be some global policies and tactics that should be discussed here. Still, decisions about interfaces, algorithms, and/or data structures might be more appropriately discussed in the same (sub) section as its corresponding software component in one of these other sections.
